Long-running agent harnesses accumulate three kinds of state: knowledge about the environment, procedural skills, and deontic state (what the agent may, must not, or may no longer do). Knowledge and skills go stale and should decay; deontic state does not go stale by disuse. Yet every published memory-decay mechanism scores entries by recency, access or outcome, so a once-stated revocation is the first thing pruned, and the consolidation step that shrinks memory is where recent work shows LLM writers fabricate authority. We propose the Typed Harness State Model (THSM): harness state is typed (episodic, semantic, procedural, deontic), carries an integrity label from a Biba-style lattice, and is governed by operators that apply decay and consolidation only to knowledge and skills while deontic state changes only through provenance-backed principal events and is enforced by a deterministic gate. Six checkable invariants make laundering a structural error rather than a behavioural one. We introduce a dual benchmark that scores any memory configuration on a utility axis and an authority axis over the same event-sourced scenarios, with the authority axis graded from tool calls only. Across 27 live experiments, four model families (gpt-4o-mini, gemini-2.5-flash-lite, qwen3-coder-30b, claude-sonnet-5), two domains and about 800 THSM cells, THSM produced zero false authority at utility within $\pm$0.05 of matched type-blind stores, while type-blind memory carried out 27--84\% of unauthorized requests depending on model and decay level. Three findings go beyond the headline: (i) pinning the deontic state into context every turn, perfectly preserved, still leaves 20--60\% false authority, and models restate the authority state correctly (92--95\%) while acting against it; (ii) skills carry authority: a task whose grant was revoked is still executed 39--97\% of the time from type-blind memory on every model, and 66--92\% even with the revocation pinned on the three smaller models (3\% on claude-sonnet-5); (iii) integrity labels without a trusted deontic channel make things worse, because flagging every permission note as unverified erodes the prohibitions too. We also reproduce, in a persistent-memory setting, two recent in-context results: constraint erosion with depth for ACT-R-style decay, and the limited role of compaction once constraints live in memory. Finally, we validate the utility half of the claim on Continual-ARC, a skill-retention benchmark built for a different paper: the exemption's cost is bounded by what decay actually evicts, and only power-law (ACT-R) activation evicts under realistic recurrence, which is the same policy whose false authority rises monotonically with decay in our own benchmark.
